\section{Relier la carte au texte}
\label{sec:ancrage}

\subsection{Citer n’est pas prouver}

ALCE évalue des réponses générées avec citations selon trois dimensions : la fluidité, l’exactitude et la qualité des citations, cette dernière mesurée par un modèle d’inférence textuelle qui vérifie si les passages cités impliquent ce qui est affirmé \parencite[6466]{ref17}. Sur le jeu de données ELI5, même les meilleurs modèles n’ont pas de citations qui soutiennent entièrement leur réponse dans environ la moitié des cas \parencite[6465]{ref17}. Une citation bien placée ne garantit donc pas le soutien. SciFact pose la question dans les termes les plus proches des nôtres : décider si des résumés d’articles soutiennent ou réfutent une affirmation scientifique, et désigner les phrases qui justifient cette décision \parencite[7534]{ref28}.

Le prototype sépare ainsi deux contrôles. Le premier est mécanique : l’extrait cité doit se retrouver dans la phrase désignée par son identifiant, par comparaison approximative avec un seuil réglable ; si l’identifiant ne convient pas, l’extrait est recherché dans tout le texte. Le second est un jugement : un appel distinct du modèle décide si les phrases citées soutiennent l’étape ou la relation, entièrement, en partie ou pas du tout. Chaque élément reçoit un statut, \enquote{vérifié}, \enquote{appui partiel} ou \enquote{à vérifier}, qui s’affiche sur la carte ; ce qui reste douteux après deux tours de correction apparaît dans un cadre en pointillés. Story Ribbons pratique le même contrôle littéral pour ses citations \parencite[3]{ref01}.

\subsection{Lire dans le document}

Le projet Semantic Reader \parencite{ref15} a exploré, à travers dix prototypes étudiés avec plus de 300 participants, comment augmenter un article sans le remplacer, y compris au-dessus des PDF existants \parencite[1]{ref15}. ScholarPhi affiche la définition des termes et des symboles dans des infobulles, Scim surligne les passages selon leur fonction dans le discours, Paper Plain propose des résumés en langage simple de certains passages \parencite[4--6]{ref15}. Les fiches de termes du prototype et le surlignage des passages cités dans le texte intégral relèvent de cette famille. SciDaSynth \parencite{ref16} applique le même principe à l’extraction de données : chaque valeur extraite par le LLM reste liée à sa source, pour être validée, corrigée ou affinée. Dans une étude intra-sujets avec 24 chercheurs, les participants obtiennent des données de bonne qualité en nettement moins de temps qu’avec les méthodes de référence \parencite[2, 10]{ref16}.

Voir le passage dans sa mise en page d’origine aide le lecteur à le reconnaître et à le situer dans la page, près d’une figure ou dans une colonne. C’est pourquoi le prototype affiche l’extrait rendu depuis le PDF plutôt qu’une simple chaîne de caractères. Cela aide à juger une étape ; cela ne prouve pas la relation proposée.

\subsection{Accéder au texte}

Tout ce qui précède dépend de la lecture du document. GROBID convertit les PDF scientifiques en XML/TEI (en-tête, références, corps du texte) à l’aide de champs conditionnels aléatoires ou, en option, de réseaux profonds \parencite{ref22}. Docling s’appuie sur un modèle d’analyse de mise en page, entraîné notamment sur DocLayNet, et sur le modèle TableFormer pour les tableaux ; il reconstitue l’ordre de lecture \parencite[2--3]{ref23}. Le prototype utilise PyMuPDF, qui fournit les coordonnées de chaque mot : elles sont indispensables pour surligner une phrase dans la page. GROBID peut s’y ajouter pour enrichir la structure.

Trois sources d’erreur se cumulent : l’ordre de lecture, le découpage en phrases et l’interprétation. Une carte bien dessinée ne corrige aucune d’elles. Le premier essai du prototype sur un article de 14 pages en deux colonnes l’a montré : faute de règles de mise en page, des graduations de graphiques et des fragments de formules étaient pris pour des titres de section, et les deux colonnes se mélangeaient. Le prototype reconstitue désormais l’ordre des colonnes, repère les titres d’après la taille et la graisse des caractères, écarte figures, tableaux et formules, et recolle les phrases coupées par un changement de colonne ou de page. Les tableaux sans filets, les mises en page à trois colonnes et les PDF numérisés restent hors de portée ; un texte trop peu extractible déclenche un message explicite.
